\subsection*{日志、提交与确定性应用}

G1 由三个独立副本构成，每个副本保存自己的持久日志、快照和状态机状态。leader 接收 b1 的受理命令，先追加到日志，再复制给 follower。达到 Raft 的提交条件后，命令进入已经承诺的日志前缀，各副本按相同次序将其应用到业务状态机。G2 和 G3 分别拥有自己的多数，b1 的一次提交不会占用另外两个组的业务日志~\cite{raft}。

受理命令携带完整 $S$、规范化原输入、输入版本和本次准入判断使用的时间上界。应用过程读取会话状态，比较同 $S$ 的原输入，分配新的分区内序号，保存受理和待办。所有副本应用这条命令时，必须得到相同的状态变化与原响应。日志确定次序，状态机定义这一步在业务上做什么。

工作者的随机退避、当前墙上时间、网络响应和平台付款都在状态机外发生。若每个副本在 Apply 中读取自己的时钟，Alice 的会话可能在一个副本上有效，在另一个副本上过期。若三个副本各自在 Apply 中发起付款，一条日志命令就会展开成三个独立外部调用，重放日志还会再次调用。状态机保存付款意图和已核验的外部证据，工作者在状态机外推进这个意图。

\CSeventeenFigure{ch17-03-raft-apply}{命令先进入提交前缀，再由状态机形成实际受理和响应}{fig:ch17-raft-apply}

追加成功、提交成功和应用完成是三个不同进度。leader 的本地持久日志已经出现一条命令时，其他副本还可能没有收到；条目进入提交前缀后，业务状态机也可能尚未执行到该位置。受理接口等待命令应用完成，再返回业务状态中的实际 P/C/T。连接若在应用完成后断开，调用方不知道结果，恢复请求便用原 $S$ 读取同一受理~\cite{hashicorp-raft-173}。

日志索引描述复制层进度，C 标识业务对象。leader 更替可能带来内部空条目或配置条目，两个成功购买也不必占据相邻的 Raft 索引。业务编号由分区状态机自己的分配规则产生，日志索引指出这次状态变化位于复制历史中的什么位置。

\subsection*{可以回答读请求的副本}

查询 b1 的完整状态需要一条可说明的读取路径。旧 leader 与多数失联后，本地还拥有一份看起来完整的状态。若它直接回答 Alice 没有找到这个 $S$，而新 leader 已提交该操作，调用方就会得到落后的判断。查询接口在使用当前 leader 的状态前，要确认当前领导资格，并等待本地状态机应用到读取所要求的进度。

示例权威读取采用库提供的 leader 核验和应用屏障：先让当前节点确认自己能够联系到所需多数，再把屏障之前的工作应用完毕，随后在状态锁下取得一致图像。任一步返回领导权丢失、超时或关闭，接口就结束本次权威读取，由调用方定位当前 leader 后重试。这条路径把领导资格、应用前缀和业务读取连接起来，更换库或读取优化时也要逐项核对这些保证~\cite{hashicorp-raft-173}。

\CSeventeenFigure{ch17-17-authoritative-read}{权威读取通过领导资格核验和应用屏障后取得状态图像}{fig:ch17-authoritative-read}

对于只展示近似进度的界面，系统可以另设允许陈旧的查询投影。投影保存来源和水位，入口据此解释付款事实与展示进度之间的差异。首次受理的唯一性判断和退款占额判断在有资格的业务状态机中进行，派生副本上的缺席只说明这份投影还没看见相应事实。

\subsection*{快照与日志截断}

随着日志增长，副本把已应用的完整状态保存成快照，并在快照可靠保存后回收其覆盖的部分日志。快照对应一个确定的应用位置，会话、绑定、原结果、付款待办、退款占额、发布标记、消费交接回执及进度和事务责任都属于这个位置上的状态。快照之后的修改由后续日志补上。

快照对象与持续应用之间也有并发边界。实现可以在状态锁下复制一份不可变图像，再在锁外写入快照存储。复制完成后，新受理可以继续应用，正在写出的图像保持原值。若写出过程直接遍历还在变化的业务映射，文件可能混入较新的订单和较旧的退款占额，最终无法对应任何实际应用前缀。

\CSeventeenFigure{ch17-04-snapshot}{完整快照覆盖一个应用前缀，后续日志补齐恢复位置}{fig:ch17-snapshot}

快照封装保存格式版本、应用位置、完整状态、分对象计数和状态摘要。恢复先解析并核验封装，再检查记录间的联系，最后一次性替换当前状态。缺失 \code{operations}、缺失退款映射或把 C1 改成新编号，均在替换前拒绝。恢复失败时原在线状态保持原值，运维可继续定位失败图像。

摘要能够发现字节损坏和与清单不符的遗漏，前提是清单及摘要来自可信保存路径。操作者若重新制作一份删掉旧购买的图像，又同时改写清单和摘要，普通哈希无法证明历史真实性。一份完整但过旧的快照也可以通过内部校验，却缺少平台上已经发生的付款。灾难恢复根据独立保存的业务历史决定开放范围，内部完整性校验负责判断取得的图像是否符合其封装与结构。

三副本多数复制依赖各副本自己的持久存储和故障域。三个进程共用一块失效磁盘，无法承受该磁盘同时丢失；三个副本分散在三个地域，则把多数通信放进跨地域路径。组内复制的部署位置、网络和保存条件，直接进入第二十一章的延迟与恢复预算\footnote{本章采用 HashiCorp Raft v1.7.3 与 raft-boltdb/v2 v2.3.1。应用屏障及快照接口按该版本的 \code{api.go}、\code{fsm.go} 核对；具体运行配置保存在配套材料。}。

\subsection*{从多数复制换到链式复制}

同一份业务状态也可以放进另一种复制结构。Chain Replication 把副本排成有方向的链：更新从链头进入，沿链传播，查询由链尾回答，更新完成也以链尾为准。链头安排更新次序并计算更新后的值，下游接收这次更新产生的结果或状态变化。原协议允许链头把一次非确定选择的结果向后传递，下游据此保存相同结果~\cite{chain-replication2004}。

设链为 $\mathrm{H}\to\mathrm{M}\to\mathrm{T}$，三者当前都保存到更新 41。更新 42 先在 H 持久保存，再传给 M，最后到达 T。H、M、T 的已处理历史分别构成由长到短的前缀关系。T 接收 42 后，链尾可把 42 作为已完成更新响应给调用方，并按顺序处理后来的查询。查询若随意落到 H 或 M，就可能越过尚未到达链尾的完成边界。

\CSeventeenFigure{ch17-05-chain}{链头排列更新，链尾形成完成边界并回答查询}{fig:ch17-chain}

每个中间副本保存 Sent 集合，记录已向后传播、尚未获得链尾完成确认的更新。链尾确认沿反方向返回，上游据此删除相应 Sent。确认丢失后，上游还可能保留已经完成的更新，恢复时重新发送相同序号和内容，下游按原传播序号核对。传播序号服务于链内修复，长期业务身份 $S$ 和原结果由被复制的应用状态保存。

现在让 M 在持久保存 42 后故障，T 还只有 41。配置管理者先与后继 T 确定其已接收序号，再通知前驱 H 改为连接 T。H 根据 Sent 将缺失的 42 补给 T，接好这段后缀后才放行后来的更新。管理者若只改 next 指针，H 接下来的 43 可能先到 T，链尾就会面对缺口。

\CSeventeenFigure{ch17-06-chain-repair}{中间节点失效后，前驱按后继水位补齐 Sent 后缀}{fig:ch17-chain-repair}

头节点故障和尾节点故障有不同的恢复结果。若 42 只到达 H，H 在向下游传播前故障，客户端尚未得到完成响应，新链头可以没有 42。客户端按原 $S$ 恢复，由新的业务状态决定首次受理或返回原结果。若 T 故障，M 已保存的在途更新可能在 M 成为新尾时变成完成更新；此前已由旧 T 回复的前缀则全部保留。

增加新尾也要交接服务位置。旧尾先向新尾发送状态图像，复制期间记录后来的增量，再停止旧尾的对外尾服务，把最后后缀补齐，配置管理者才发布新尾位置。图像、增量与服务资格依此连接，查询和完成响应便沿同一条历史前进。

\CSeventeenFigure{ch17-18-new-tail}{增加新尾先复制图像与增量，再交接尾部服务资格}{fig:ch17-new-tail}

这些动作依赖有序可靠链路、失效停止的副本和有资格的配置管理者~\cite{chain-replication2004}\footnote{原论文的故障检测与配置服务承担这一前提。配套 TCP 机制对照采用可信串行管理者，移除进程前先确认它已经停止，未实现自动网络分区重配置。}。实际系统要复制并保护管理状态。网络分区中，单节点的超时观察不足以证明另一端停止服务，重配置必须先阻止旧服务与新链并行对外确认不同历史。
\FloatBarrier
